% !PRINTED 518
\weberpage{518}

%<<CONT
certainly a group. If $\pi_\beta$ is a permutation by which the
row $A$ passes into the row $B$, which by the hypothesis of
transitivity must exist in $P$, then
$Q_\alpha \pi_\beta$ is a subsidiary group of $Q_\alpha$ by which
$A$ passes into $B$. $\pi_\beta^{-1} Q_\alpha \pi_\beta =
Q_\beta$ is a group conjugate to $Q_\alpha$, and indeed the one
by whose permutations the row $B$ is not displaced.
In this way we obtain the conjugate groups

\[ Q_\alpha,\ Q_\beta \dots Q_\sigma, \]

whose greatest common divisor is the group $Q$ of index $j$
considered above.

A function $y_\alpha$ belonging to $Q_\alpha$ is the root
of an equation of degree $s$ in $\Omega$,
$\varphi(y) = 0$, and the remaining roots $y_\beta \dots y_\sigma$
arise from $y_\alpha$ by application of the permutations
of the subsidiary groups $Q_\alpha \pi_\beta, \dots Q_\alpha \pi_\sigma$.
The equation $\chi(u) = 0$ of degree $j$ found above is
the Galois resolvent of the equation $\varphi(y) = 0$. We
have remarked above that the group of the equation $\chi(u) = 0$
of degree $j$ is isomorphic with the group of the permutations
that are produced by $P$ among the systems of the $A, B \dots S$.
But since, on account of the transitivity of $P$,
these systems also pass transitively into one another, $j$
must, by the theorem of \S.\,154, 7., be divisible by $s$, and in the
particular case where $j$ is a prime, we must have $j = s$.

The simultaneous adjunction of all the roots of $\varphi(y)$
has the same effect as the adjunction of the function $\psi$,
namely, to reduce the group of the equation $f(x) = 0$ from $P$ to $Q$.

By the adjunction of $y_\alpha$ a rational factor
$f(x, y_\alpha)$ splits off from $f(x)$, whose roots are the
quantities of the row $A$, and by adjunction of the remaining $y$
one obtains the decomposition

\weq{6}{f(x) = f(x,y_\alpha)\, f(x,y_\beta) \dots f(x,y_\sigma),}

The functions $f(x,y_\alpha), f(x,y_\beta) \dots f(x,y_\sigma)$
are precisely the factors $f_\alpha(x,\psi), f_\beta(x,\psi) \dots
f_\sigma(x,\psi)$ in (2).

It is further to be remarked here that $f(x,y_\alpha)$ is
irreducible in the field $\Omega(y_\alpha)$. For on account of the
assumed irreducibility of $f(x)$ there are in $P$
permutations by which $\alpha$ is carried into every $\alpha_i$,
and these permutations belong to
$Q_\alpha$. If therefore a rational function $\varphi(x,
y_\alpha)$ vanishes for $x = \alpha$,
%CONT>>
